\section{Filtered rings I}
\begin{definition}[$\mathbb Z$-filtration]\label{def:z_filtration}
    Let $A$ be a unital ring. \(\Gamma_\bullet A\) is a family \(\{\Gamma_i A\}_{i\in\mathbb Z}\) called a \textbf{$\mathbb Z$-filtration} of $A$ if
    \begin{itemize}
        \item Each \(\Gamma_i A\) is an additive subgroup of \(A\);
        \item $\Gamma_i A \subseteq A$;
        \item $\Gamma_i A \subseteq \Gamma_{i+1} A$ and $1\in \Gamma_0 A$;
        \item $\bigcup_{i\in\mathbb Z} \Gamma_i A = A$;
        \item $\Gamma_i A\cdot \Gamma_j A \subseteq \Gamma_{i+j} A$ for all $i,j$; and often we assume the
              filtration is separated: $\bigcap_{i\in\mathbb Z}\Gamma_i A=0$.
    \end{itemize}
\end{definition}
We say $A$ is \textbf{$\mathbb Z$-filtered} if $\Gamma_\bullet A$ is$\mathbb Z_{\ge 0}$-filtered if additionally $\Gamma_i A=0$ for all $i<0$.

The \textbf{associated graded ring} is
\[
    \operatorname{gr}_\bullet A = \bigoplus_{i\in\mathbb Z} \Gamma_i A/\Gamma_{i-1} A.
\]
The Rees ring (with respect to $\Gamma$) is
\[
    R^{\Gamma}_\bullet(A) = \bigoplus_{i\in\mathbb Z} \Gamma_i A\,u^i \subseteq A[u,u^{-1}],
\]
which records the filtration multiplicatively.

\subsection*{Basic consequences}
\begin{lemma}\label{lem:rees-quotients}
    For the Rees ring $R^{\Gamma}_\bullet (A)$ one has natural isomorphisms
    \[
        \gr_\bullet (A) \simeq R^{\Gamma}_\bullet (A)/u\,R^{\Gamma}_\bullet (A),\qquad
        A \simeq R^{\Gamma}_\bullet (A)/(u-1)\,R^{\Gamma}_\bullet (A).
    \]
\end{lemma}

\begin{proposition}\label{prop:noeth}
    Let $A$ be a filtered ring with filtration $\Gamma$.
    \begin{enumerate}
        \item If $R^{\Gamma}_\bullet (A)$ is Noetherian, then both $A$ and $\gr_\bullet (A)$ are Noetherian.
        \item If the filtration is nonnegative (i.e. $\Gamma_iA=0$ for $i<0$) and
              $\gr_\bullet (A)$ is Noetherian, then $A$ is Noetherian.
    \end{enumerate}
\end{proposition}
\begin{proof}
    (1) By Lemma~\ref{lem:rees-quotients}, both $\operatorname{gr}_\bullet (A)$ and $A$ are surjective homomorphic images of the Rees ring $R^{\Gamma}_\bullet (A)$. Since the homomorphic image of a Noetherian ring is Noetherian, if $R^{\Gamma}_\bullet (A)$ is Noetherian, then so are its quotients $\operatorname{gr}_\bullet (A)$ and $A$.

    (2) Assume the filtration is nonnegative, i.e., $\Gamma_i A = 0$ for $i < 0$, which implies $A = \bigcup_{p \ge 0} \Gamma_p A$.
    Let $I_1 \subseteq I_2 \subseteq \cdots$ be an ascending chain of left ideals of $A$.
    Consider the associated graded ideals in $\operatorname{gr}_\bullet (A)$ defined by
    \[
        \mathfrak{a}_j := \operatorname{gr}_\bullet(I_j) = \bigoplus_{p \ge 0} \frac{I_j \cap \Gamma_p A + \Gamma_{p-1} A}{\Gamma_{p-1} A}.
    \]
    Since $I_j \subseteq I_{j+1}$, we have an ascending chain $\mathfrak{a}_1 \subseteq \mathfrak{a}_2 \subseteq \cdots$ in $\operatorname{gr}_\bullet (A)$.
    Because $\operatorname{gr}_\bullet (A)$ is Noetherian, this chain stabilizes. Choose $N$ such that $\operatorname{gr}_\bullet(I_j) = \operatorname{gr}_\bullet(I_N)$ for all $j \ge N$.

    We claim that $I_j = I_N$ for all $j \ge N$, which implies $A$ is Noetherian.
    Fix $j \ge N$. Since $I_N \subseteq I_j$, it suffices to show equality. We proceed by induction on the filtration degree $p$ that
    \[
        I_j \cap \Gamma_p A = I_N \cap \Gamma_p A.
    \]
    For $p < 0$, both sides are $0$. Assume the equality holds for $p-1$.
    Since $\operatorname{gr}_\bullet(I_j) = \operatorname{gr}_\bullet(I_N)$, their degree-$p$ components are identical:
    \[
        \frac{I_j \cap \Gamma_p A + \Gamma_{p-1} A}{\Gamma_{p-1} A} = \frac{I_N \cap \Gamma_p A + \Gamma_{p-1} A}{\Gamma_{p-1} A}.
    \]
    This equality implies
    \[
        I_j \cap \Gamma_p A + \Gamma_{p-1} A = I_N \cap \Gamma_p A + \Gamma_{p-1} A.
    \]
    Intersecting both sides with $I_j \cap \Gamma_p A$ and using the modular law ($I_N \subseteq I_j$, so the modular law applies), we get:
    \[
        I_j \cap \Gamma_p A = I_N \cap \Gamma_p A + (I_j \cap \Gamma_{p-1} A).
    \]
    By the induction hypothesis, $I_j \cap \Gamma_{p-1} A = I_N \cap \Gamma_{p-1} A \subseteq I_N \cap \Gamma_p A$. Thus,
    \[
        I_j \cap \Gamma_p A = I_N \cap \Gamma_p A.
    \]
    Since $A = \bigcup_{p} \Gamma_p A$, taking the union over all $p$ yields $I_j = I_N$.
    Thus, every ascending chain stabilizes, and $A$ is Noetherian.
\end{proof}

\begin{remark}
    It is not recommended to use ``s.t.'' to replace ``such that'' in the formal writing,
    since ``s.t.'' may represent ``subject to'' in some contexts.
\end{remark}

\subsection*{Filtered modules and good filtrations}
Let $A$ be a $\mathbb{Z}$-filtered ring with filtration $\Gamma_\bullet A$ and let $M$ be a left
$A$-module.

\begin{definition}[Compatible and Good Filtration]\label{def:filtered-module}
    Let $\Omega_\bullet M$ be a $\mathbb{Z}$-filtration on $M$ (also, each $\Omega_i M$ is an additive subgroup of $M$) such that $\bigcup_k \Omega_k M = M$.

    \begin{enumerate}
        \item $\Omega_\bullet M$ is called a \textbf{compatible filtration} over $\Gamma_\bullet A$ if:
              \[
                  \Gamma_k A \cdot \Omega_i M \subseteq \Omega_{k+i} M.
              \]

        \item $\Omega_\bullet M$ is called a \textbf{good filtration} over $\Gamma_\bullet A$ if it is compatible and the Rees module:
              \[
                  R_\bullet(M) = R_\bullet^\Omega(M) = \bigoplus_i \Omega_i M \cdot u^i
              \]
              is finitely generated (f.g.) over $R(A)$.

              This implies that the associated graded module:
              \[
                  \gr_\bullet M = \gr_\bullet^\Omega M = \bigoplus_i \frac{\Omega_i M}{\Omega_{i-1} M}
              \]
              is finitely generated (f.g.) over $\gr_\bullet A$.
    \end{enumerate}
\end{definition}

\begin{lemma}\label{lem:good-filtration-exists}
    \mbox{}
    \begin{enumerate}
        \item If $M$ is finitely generated over $A$, then good filtrations on $M$ exist.
        \item If $\Omega_\bullet M$ is good, then there exist integers $k_1,\dots,k_\ell$ and
              elements $m_1,\dots,m_\ell\in M$ with $m_i\in\Omega_{k_i}M$ such that, for every $p$,
              \[
                  \Omega_pM \,=\, \sum_{i=1}^{\ell} \Gamma_{p-k_i}A\cdot m_i.
              \]
    \end{enumerate}
\end{lemma}
\begin{proof}
    \textbf{(1)} Assume $M$ is generated by $m_1, \dots, m_\ell$ over $A$.
    Then define the filtration by:
    \[
        \FilM{p} = \sum_{i=1}^{\ell} \FilA{p} \cdot m_i.
    \]
    This defines a good filtration on $M$.

    \vspace{1em}

    \textbf{(2)} By definition, if $\FilM{\bullet}$ is a good filtration, then the Rees module $\Rees^\Omega(M)$ is finitely generated (f.g.) over the Rees ring $\Rees(A)$.

    This implies we can find homogeneous generators:
    \[
        m_1 u^{k_1}, \dots, m_\ell u^{k_\ell}
    \]
    where $m_i \in \FilM{k_i}$ (considered as degree $k_i$ elements).

    Expressing the degree $p$ part of the Rees module using these generators:
    \[
        \FilM{p} \cdot u^p = \sum_{i=1}^{\ell} \left( \FilA{p-k_i} \cdot u^{p-k_i} \right) \cdot (m_i \cdot u^{k_i}).
    \]
    By canceling $u^p$ on both sides (or comparing coefficients), we obtain the desired formula:
    \[
        \FilM{p} = \sum_{i=1}^{\ell} \FilA{p-k_i} \cdot m_i.
    \]

\end{proof}

\begin{corollary}\label{cor:compare-good-filtrations}
    Let $M$ be a finitely generated $A$-module and let $\Omega_\bullet M$ and $F_\bullet M$
    be two compatible filtrations on $M$.
    If $F_\bullet M$ is good, then there exists an integer $a$ such that
    \[
        F_pM \subseteq \Omega_{p+a}M \qquad \text{for all } p.
    \]
    Consequently, if $\Omega_\bullet M$ is also good, then there exists an integer $a$ with
    \[
        \Omega_{p-a}M \subseteq F_pM \subseteq \Omega_{p+a}M \qquad \text{for all } p.
    \]
\end{corollary}
\begin{proof}
    By Lemma~\ref{lem:good-filtration-exists}(2), since $\FilF{\bullet}$ is a good filtration, for all $p$ we have:
    \[
        \FilF{p} = \sum_{i=1}^{\ell} \FilA{p-k_i} \cdot m_i
    \]
    for some generators $m_1, \dots, m_\ell$ and integers $k_1, \dots, k_\ell$.

    Assume that $m_1, \dots, m_\ell \in \FilM{q}$ for some sufficiently large integer $q \gg 0$.

    Then we have the following inclusion chain:
    \begin{align*}
        \FilF{p} & \subseteq \sum_{i=1}^{\ell} \FilA{p-k_i} \cdot \FilM{q}                                                                    \\
                 & \subseteq \sum_{i=1}^{\ell} \FilM{p+q-k_i} \quad (\text{since } \FilM{\bullet} \text{ is compatible with } \FilA{\bullet}) \\
                 & \subseteq \FilM{p+a}
    \end{align*}
    where we choose $a \ge \max_i(q - k_i)$.
\end{proof}


\paragraph{Grothendieck group.}
For any ring $A$, let $K_0(A)$ denote the Grothendieck group of the abelian
category $\Mod(A)$: it is generated by symbols $[M]$ for $M\in\Mod(A)$ with the
relations $[M]=[M_1]+[M_2]$ for every short exact sequence
$0\to M_1\to M\to M_2\to 0$.

\begin{remark}[Well-definedness of $K_0$]
    Strictly speaking, defining the Grothendieck group on the category of \textit{all} modules $\mathrm{Mod}(A)$ is problematic for two reasons:

    \begin{enumerate}
        \item \textbf{Set-theoretic Size Issue}: The collection of isomorphism classes in $\mathrm{Mod}(A)$ forms a \textbf{proper class}, not a set. One cannot strictly generate a free abelian group from a proper class in ZFC.

        \item \textbf{Eilenberg Swindle}: Allowing infinite direct sums makes the group trivial. For any module $M$, let $M^\infty = \bigoplus_{i=1}^\infty M$. We have the isomorphism:
              \[
                  M \oplus M^\infty \cong M^\infty.
              \]
              In the Grothendieck group, this implies $[M] + [M^\infty] = [M^\infty]$, and by cancellation, $[M] = 0$.
    \end{enumerate}

    \textbf{Correction}: To avoid these issues, $K$-theory is standardly defined on "small" subcategories of \textbf{finitely generated} modules (where infinite sums are not allowed and isomorphism classes form a set):
    \begin{itemize}
        \item $K_0(A)$ is defined on $\mathbf{Proj}(A)$ (f.g. projective modules).
        \item $G_0(A)$ is defined on $\mathbf{mod}(A)$ (f.g. modules, typically for Noetherian rings).
    \end{itemize}
\end{remark}
